\begin{table}[H]
\centering
\caption{Robustness Checks}
\begin{threeparttable}
\begin{tabular}{lcc}
\toprule
 & Coefficient & SE \\
\midrule
\multicolumn{3}{l}{\textit{Panel A: Baseline (Table 3, Col. 1)}} \\
\quad log(Initial claims) & 10.96 & (12.62) \\
 & & \\
\multicolumn{3}{l}{\textit{Panel B: Leave-one-industry-out}} \\
\quad Drop 23 & 11.87 & (13.62) \\
\quad Drop 31-33 & 20.89 & (22.86) \\
\quad Drop 56 & 11.80 & (12.18) \\
\quad Drop 44-45 & 11.09 & (12.51) \\
\quad Drop 21 & 5.97 & (12.47) \\
 & & \\
\multicolumn{3}{l}{\textit{Panel C: Alternative outcome thresholds (2SLS)}} \\
\quad 7-day timeliness & 14.29 & (14.56) \\
\quad 21-day timeliness & 7.18 & (8.16) \\
 & & \\
\multicolumn{3}{l}{\textit{Panel D: Placebo (pre-recession, 2006--2007)}} \\
\quad Bartik shock (reduced form) & 593.5 & (426.4) \\
\bottomrule
\end{tabular}
\begin{tablenotes}[flushleft]
\small
\item \textit{Notes:} Panel A repeats the baseline 2SLS estimate. Panel B drops the highest-Rotemberg-weight industry and re-estimates. Panel C uses alternative timeliness thresholds (7-day, 21-day). Panel D runs the reduced form on pre-recession years only (placebo test). Panel E reports wild cluster bootstrap confidence intervals. All specifications include state and year fixed effects with standard errors clustered at the state level.
\end{tablenotes}
\end{threeparttable}
\label{tab:robust}
\end{table}
